\newpage

\section{Анализ предикторов успеха}

В данном разделе проводится сравнительный анализ групп \textbf{success} (767 компаний) и \textbf{failure} (1011 компаний). Группа active (3961 компания) исключена из анализа, так как итоговый исход для этих компаний еще не известен.

Во избежание ошибок при анализе в данной работе применяется следующие подходы:

\begin{itemize}
    \item \textbf{Контроль по когортам:} для каждой гипотезы приводится не только общее сравнение, но и анализ по году батча, что позволяет сравнивать компании одного <<возраста>>.
    \item \textbf{Фокус на признаках, известных на момент входа в акселератор:} тип команды основателей, отрасль, сезон батча, лаконичность описания, количество тегов, география.
\end{itemize}

Для проверки статистической значимости различий между группами используется критерий Манна-Уитни~\cite{mann-whitney}.

\subsection{Гипотеза 1: Соло-основатели vs команды}

\textbf{H1:} стартапы с единственным основателем имеют более низкую долю успешных исходов по сравнению с командами из нескольких основателей.

\textbf{Обоснование:} единственный основатель с высокой долей вероятности будет иметь более узкие компетенции, а также рискует однобоко мыслить при принятии решений на ранних этапах.

Анализ (Рис.~\ref{fig:h1}) показал выраженную зависимость: доля успешных стартапов монотонно растет с увеличением числа основателей до трёх и более. Основатели-одиночки (Solo) имеют долю успешных 35.1\% ($n = 518$), пары (Pair) --- 45.7\% ($n = 831$), команды из трех и более человек (Team 3+) --- 50.9\% ($n = 381$).

\begin{figure}[htbp]
    \centering
    \includegraphics[width=0.95\textwidth]{h1_solo_vs_cofounders.png}
    \caption{Гипотеза 1: Соло-основатели vs команды основателей}
    \label{fig:h1}
\end{figure}

При этом когортный анализ (правая панель Рис.~\ref{fig:h1}) показывает, что доли успешных стартапов по числу основателей хотя и изменяются при контроле по году батча, однако среди компаний с тремя и более основателями доля успешных всегда либо превосходит доли успешных стартапов с одним или двумя основателями, либо лишь незначительно уступает.

Критерий Манна-Уитни подтвердил статистическую значимость различия между соло-основателями и командами ($p < 0.05$).

\textbf{Вывод:} гипотеза подтверждена. Наличие трех и более основателей положительно связано с вероятностью успеха. Разница между соло-основателями и командами составляет около 15 процентных пунктов.

\subsection{Гипотеза 2: B2B vs Consumer}

\textbf{H2:} B2B-стартапы имеют более высокую долю успешных по сравнению с Consumer-стартапами.

\textbf{Обоснование:} B2B-компании имеют более ясную решаемую проблему, предсказуемую модель монетизации и чаще становятся целями M\&A.

Результаты (Рис.~\ref{fig:h2}) показали существенное различие: доля успешных среди B2B-стартапов составляет 49.1\% ($n = 851$), тогда как среди Consumer --- лишь 32.9\% ($n = 459$). Разница составляет 16 процентных пунктов.

\begin{figure}[htbp]
    \centering
    \includegraphics[width=0.95\textwidth]{h2_b2b_vs_consumer.png}
    \caption{Гипотеза 2: B2B vs Consumer стартапы}
    \label{fig:h2}
\end{figure}

Когортный анализ демонстрирует устойчивость этого различия: линия B2B стабильно расположена выше линии Consumer почти во всех когортах 2005--2025 годов.

Критерий Манна-Уитни подтвердил статистическую значимость различия ($p < 0.05$).

\textbf{Вывод:} гипотеза подтверждена. B2B-стартапы имеют значимо более высокую долю успешных, и это различие устойчиво при контроле по когортам.

\subsection{Гипотеза 3: Сезон батча}

\textbf{H3:} зимние батчи (Winter) имеют иную долю успешных, чем летние (Summer).

\textbf{Обоснование:} различия в конкурсе, рыночных условиях и составе участников могут влиять на исход.

Анализ (Рис.~\ref{fig:h3}) показал минимальное различие: доля успешных среди летних батчей составляет 42.5\% ($n = 872$), среди зимних --- 44.1\% ($n = 923$). Разница менее 2 процентных пунктов.

\begin{figure}[htbp]
    \centering
    \includegraphics[width=0.95\textwidth]{h3_batch_season.png}
    \caption{Гипотеза 3: Сезон батча (Summer vs Winter)}
    \label{fig:h3}
\end{figure}

Когортный анализ подтверждает отсутствие устойчивого паттерна: линии Summer и Winter многократно пересекаются, не демонстрируя систематического преимущества одного сезона над другим.

\textbf{Вывод:} гипотеза не подтверждена. Сезон батча не является значимым предиктором успеха стартапа, то есть время входа в акселератор (зима или лето) не влияет на итоговый исход.

\subsection{Гипотеза 4: Лаконичность описания}

\textbf{H4:} стартапы с более коротким описанием чаще достигают успеха.

\textbf{Обоснование:} лаконичное описание может свидетельствовать о ясности бизнес-модели и четком понимании продукта.

Анализ (Рис.~\ref{fig:h4}) выявил нелинейную зависимость. Boxplot показывает, что медианы длины описания для групп success и failure близки (около 52 и 51 символа соответственно). Однако анализ по квартилям длины обнаружил следующую закономерность: наибольшая доля успешных наблюдается в квартилях Q2 (48.1\%, $n = 393$) и Q3 (49.0\%, $n = 455$), тогда как самые короткие (Q1: 40.0\%, $n = 445$) и самые длинные (Q4: 41.5\%, $n = 381$) описания связаны с более низкой долей успеха.

\begin{figure}[htbp]
    \centering
    \includegraphics[width=0.95\textwidth]{h4_desc_length.png}
    \caption{Гипотеза 4: Длина описания и успех стартапа}
    \label{fig:h4}
\end{figure}

\textbf{Вывод:} гипотеза частично подтверждена в модифицированном виде. Не самые короткие, а \emph{умеренно лаконичные} описания (Q2--Q3) связаны с наибольшей долей успеха. Слишком короткие описания могут свидетельствовать о недостаточной проработке идеи, а слишком длинные --- о размытости позиционирования.

\subsection{Гипотеза 5: Количество тегов}

\textbf{H5:} стартапы с большим числом тегов имеют более размытое позиционирование и как следствие чаще терпят неудачу.

\textbf{Обоснование:} фокусировка на узкой нише эффективнее, чем попытка охватить множество направлений.

Анализ (Рис.~\ref{fig:h5}) выявил нелинейную зависимость. Наибольшая доля успешных наблюдается у стартапов с 3--4 тегами (48.2\%, $n = 498$), за ними следуют стартапы с 1--2 тегами (43.0\%, $n = 969$). Стартапы без тегов (36.3\%, $n = 204$) и с 5+ тегами (37.6\%, $n = 133$) имеют наименьшую долю успешных.

\begin{figure}[htbp]
    \centering
    \includegraphics[width=0.95\textwidth]{h5_tag_count.png}
    \caption{Гипотеза 5: Количество тегов и успех стартапа}
    \label{fig:h5}
\end{figure}

\textbf{Вывод:} гипотеза частично подтверждена. Оптимальным является умеренное количество тегов (3--4), что свидетельствует о четком, но не чрезмерно узком позиционировании. Отсутствие же тегов, также как и избыток, связаны с более низкой долей успеха.


\newpage
\subsection{Гипотеза 6: Активный найм как сигнал}

\textbf{H6:} стартапы с флагом \texttt{isHiring = True} реже имеют статус inactive.

\textbf{Обоснование:} активный найм является прокси-индикатором наличия финансирования и роста компании.

Результаты (Рис.~\ref{fig:h6}) показали крайне выраженную связь: среди нанимающих компаний доля успешных составляет практически 100\%, тогда как среди не нанимающих --- около 43\%.

\begin{figure}[htbp]
    \centering
    \includegraphics[width=0.95\textwidth]{h6_hiring_signal.png}
    \caption{Гипотеза 6: Активный найм и успех стартапа}
    \label{fig:h6}
\end{figure}

Когортный анализ подтверждает: линия <<Нанимает>> стабильно находится на уровне 100\% во всех когортах.

\textbf{Вывод:} связь подтверждена, однако эту информацию следует интерпретировать с осторожностью. Флаг \texttt{isHiring} отражает \emph{текущее} состояние компании (на момент сбора данных в 2025 году), а не характеристику на момент входа в акселератор. Фактически, неактивные компании по определению не могут нанимать сотрудников. Таким образом, \texttt{isHiring} является скорее \textbf{индикатором текущей активности}, чем предиктором будущего успеха. Тем не менее, для инвесторов этот признак может служить полезным сигналом при оценке текущего состояния портфельных компаний.

\subsection{Гипотеза 7: США vs мир}

\textbf{H7:} стартапы из США имеют более высокую долю успешных по сравнению с компаниями из других стран.

\textbf{Обоснование:} близость к Кремниевой долине, доступ к венчурному капиталу и развитая экосистема.

Анализ (Рис.~\ref{fig:h7}) показал умеренное различие: доля успешных среди американских стартапов составляет 45.6\% ($n = 1348$), среди стартапов из других стран --- 41.8\% ($n = 306$). Разница составляет около 4 процентных пунктов.

\begin{figure}[htbp]
    \centering
    \includegraphics[width=0.95\textwidth]{h7_us_vs_world.png}
    \caption{Гипотеза 7: США vs другие страны}
    \label{fig:h7}
\end{figure}

Когортный анализ показывает высокую волатильность для группы <<Другие>> из-за малого размера выборки в ранних когортах. В более поздних когортах (2013--2025) различие между группами минимально.

\textbf{Вывод:} различие существует, но невелико. Преимущество американских стартапов составляет около 4 процентных пунктов и может объясняться как реальным преимуществом среды, так и особенностями выборки (Y Combinator --- американский акселератор, что может создавать смещение отбора).

\subsection{Гипотеза 8: Стадия развития (Early vs Growth)}

\textbf{H8:} стартапы, попавшие в батч на стадии Growth, имеют более высокую долю успешных по сравнению со стартапами стадии Early.

\textbf{Обоснование:} стадия Growth свидетельствует о пройденном product-market fit и масштабировании.

Результаты (Рис.~\ref{fig:h8}) показали наиболее выраженное различие среди всех гипотез: доля успешных среди Growth-стартапов составляет 75.8\% ($n = 285$), тогда как среди Early --- лишь 37.2\% ($n = 1519$). Разница составляет 38.6 процентных пунктов.

\begin{figure}[htbp]
    \centering
    \includegraphics[width=0.95\textwidth]{h8_stage.png}
    \caption{Гипотеза 8: Стадия развития (Early vs Growth)}
    \label{fig:h8}
\end{figure}

Когортный анализ подтверждает устойчивость различия: линия Growth стабильно расположена значительно выше линии Early во всех когортах 2005--2025 годов, кроме, разве что, 2020, где линия Growth почти соприкасается с Early.

\textbf{Вывод:} гипотеза подтверждена. Стадия развития является сильнейшим предиктором среди проверенных гипотез. Однако, аналогично \texttt{isHiring}, стадия может отражать уже достигнутый прогресс компании, а не начальную характеристику. Тем не менее, в отличие от возраста или размера команды, стадия является качественной характеристикой бизнес-модели и может быть определена на относительно раннем этапе инвестирования.

\newpage

\subsection{Сводная таблица результатов}

Результаты проверки всех восьми гипотез представлены в Таблице~\ref{tab:hypotheses}.

\begin{table}[htbp]
\centering
\small
\caption{Сводная таблица результатов проверки гипотез}
\label{tab:hypotheses}
\begin{tabular}{@{}c p{4.5cm} p{3.5cm} c p{2.5cm}@{}}
\toprule
\textbf{\#} & \textbf{Гипотеза} & \textbf{Ключевые числа} & \textbf{$p < 0.05$} & \textbf{Результат} \\
\midrule
H1 & Соло-основатели менее успешны & Solo 35.1\% vs Team 50.9\% & Да & Подтверждена \\
\addlinespace[2pt]
H2 & B2B успешнее Consumer & B2B 49.1\% vs Consumer 32.9\% & Да & Подтверждена \\
\addlinespace[2pt]
H3 & Сезон батча влияет на успех & Summer 42.5\% vs Winter 44.1\% & Нет & Не подтверждена \\
\addlinespace[2pt]
H4 & Короткое описание = успех & Q2--Q3 $\approx$ 48.5\% vs Q1,Q4 $\approx$ 41\% & --- & Частично \\
\addlinespace[2pt]
H5 & Много тегов = неудача & 3--4 тега: 48.2\% vs 5+: 37.6\% & --- & Частично \\
\addlinespace[2pt]
H6 & Найм связан с успехом & Нанимает $\approx$ 100\% vs нет 43\% & Да & Подтверждена* \\
\addlinespace[2pt]
H7 & США успешнее & США 45.6\% vs другие 41.8\% & --- & Слабое различие \\
\addlinespace[2pt]
H8 & Growth успешнее Early & Growth 75.8\% vs Early 37.2\% & Да & Подтверждена* \\
\bottomrule
\end{tabular}

\vspace{4pt}
\raggedright\footnotesize * С оговоркой: признак может отражать текущее состояние, а не начальную характеристику.
\end{table}
